\documentclass[11pt]{article}
\usepackage[margin=1in]{geometry}
\usepackage{amsmath,amssymb}
\usepackage{enumitem}
\usepackage[T1]{fontenc}

\title{COMP2300 Set 7.5 Practice Exam Solutions}
\author{}
\date{}

\begin{document}
\maketitle

\section*{Part A: Multiple Choice}
\begin{enumerate}[label=\textbf{A\arabic*.}]
  \item B
  \item B
  \item B
  \item B
  \item B
  \item C
  \item A
  \item B
  \item C
  \item B
\end{enumerate}

\section*{Part B: Computational Problems}
\begin{enumerate}[label=\textbf{B\arabic*.}]
  \item \textbf{Hit and miss rates}
  \begin{enumerate}[label=(\alph*)]
    \item
    \[
      \text{Miss Rate} = \frac{40}{500} = 0.08 = 8\%
    \]
    \item
    \[
      \text{Hit Rate} = 1 - 0.08 = 0.92 = 92\%
    \]
  \end{enumerate}

  \item \textbf{Address breakdown}
  \begin{enumerate}[label=(\alph*)]
    \item Block size is 8 bytes:
    \[
      \log_2 8 = 3
    \]
    So there are 3 offset bits.
    \item Four sets require:
    \[
      \log_2 4 = 2
    \]
    So there are 2 set-index bits.
    \item In a 16-bit address:
    \[
      16 - 3 - 2 = 11
    \]
    So 11 tag bits remain.
  \end{enumerate}

  \item \textbf{Direct-mapped placement}
  \begin{enumerate}[label=(\alph*)]
    \item
    \[
      0x34 = 0011\,0100_2
    \]
    Low 5 bits are:
    \[
      10100
    \]
    \item Offset is the low 3 bits:
    \[
      100_2
    \]
    \item Set index is the next 2 bits:
    \[
      10_2
    \]
    which is set 2.
  \end{enumerate}

  \item \textbf{Cache benefit calculation}
  \begin{enumerate}[label=(\alph*)]
    \item At 25\% hit rate:
    \[
      AAT = (0.25)(1) + (0.75)(4) = 0.25 + 3 = 3.25
    \]
    \item At 50\% hit rate:
    \[
      AAT = (0.50)(1) + (0.50)(4) = 0.5 + 2 = 2.5
    \]
    \item Compare with 3 cycles when disabled:
    \begin{itemize}[leftmargin=1.5em]
      \item 25\% hit rate: \(3.25 > 3\), so \textbf{not beneficial}
      \item 50\% hit rate: \(2.5 < 3\), so \textbf{beneficial}
    \end{itemize}
  \end{enumerate}

  \item \textbf{SRAM vs DRAM reasoning}
  \begin{enumerate}[label=(\alph*)]
    \item SRAM is preferred for caches because it is fast and does not need refresh, which makes it suitable for latency-critical on-chip storage.
    \item DRAM is preferred for large main memory because it has much higher density and lower cost per bit, even though it is slower and requires refresh.
  \end{enumerate}

  \item \textbf{Write-policy comparison}
  \begin{enumerate}[label=(\alph*)]
    \item \textbf{write-through vs write-back}
    \begin{itemize}[leftmargin=1.5em]
      \item write-through advantage: lower-level memory is always up to date
      \item write-through tradeoff: more write traffic
      \item write-back advantage: fewer lower-level writes on repeated updates
      \item write-back tradeoff: more complexity and dirty-line management
    \end{itemize}
    \item \textbf{write-allocate vs no-write-allocate}
    \begin{itemize}[leftmargin=1.5em]
      \item write-allocate advantage: useful if the block will be written/read again soon
      \item write-allocate tradeoff: may bring in data that is not reused
      \item no-write-allocate advantage: avoids filling cache on one-off write misses
      \item no-write-allocate tradeoff: later accesses may miss again
    \end{itemize}
  \end{enumerate}
\end{enumerate}

\section*{Part C: Past Exam Questions}
\begin{enumerate}[label=\textbf{C\arabic*.}]
  \item \textbf{Cache enable decision}
  \\[0.5ex]
  Let hit rate be \(h\). Then:
  \[
    AAT = h(1) + (1-h)(4) = 4 - 3h
  \]
  Cache is beneficial when:
  \[
    4 - 3h < 3
  \]
  \[
    -3h < -1
  \]
  \[
    h > \frac{1}{3}
  \]
  So the cache is beneficial when hit rate is greater than \(33.\overline{3}\%\).

  \item \textbf{Levels in the memory hierarchy}
  \\[0.5ex]
  One reasonable ordering is:
  \[
    \text{Registers} \rightarrow \text{L1} \rightarrow \text{L2} \rightarrow \text{L3} \rightarrow \text{DRAM} \rightarrow \text{Disk/Flash}
  \]
  Lower levels are larger but slower because technology optimized for density and low cost per bit generally has longer access latency and is further from the CPU.

  \item \textbf{Cache misses in pipelined execution}
  \\[0.5ex]
  Long-latency cache misses create large stalls if the processor can only wait in order. Later out-of-order designs try to overlap other independent instructions and even multiple memory requests while a miss is being serviced, which is why memory-level parallelism becomes valuable once cache-miss latency is significant.
\end{enumerate}

\end{document}
